% from ega4-18-fr.tex
\oldpage[IV-4]{133}

Si $F=GH$ et si $G$ et $H$ engendrent l'idéal $A[T]$, il y a deux polynômes $P$, $Q$ de $A[T]$ tels que $1=PG+QH$. On en déduit que l'intersection des idéaux principaux $\mathfrak a=G.A[T]$ et $\mathfrak b=H.A[T]$ est égale à $\mathfrak c=F.A[T]$: en effet, si $R\in G.A[T]\cap H.A[T]$, on peut écrire $R=PRG+QRH$; or $RH$ (resp. $RG$) est un multiple de $F$ puisque $R$ est un multiple de $G$ (resp. $H$), donc $R\in F.A[T]$. Comme $A[T]=\mathfrak a+\mathfrak b$, $A[T]/\mathfrak c$ est somme directe des idéaux $\mathfrak a/(\mathfrak a\cap\mathfrak b)$ et $\mathfrak b/(\mathfrak a\cap\mathfrak b)$, canoniquement isomorphes respectivement à $A[T]/\mathfrak b$ et $A[T]/\mathfrak a$.

Inversement, supposons donnée une décomposition de $B$ en composée directe de deux $A$-algèbres $B'$, $B''$, qui s'identifient canoniquement à deux idéaux $e'B$, $e''B$ de $B$, correspondant à une décomposition $1=e'+e''$ de $1$ en idempotents orthogonaux $e'$, $e''$ de $B$. Posons encore $B_0=B\otimes_A k=k[T]/F_0.k[T]$, où $F_0$ est l'image canonique de $F$ dans $k[T]$, de même degré $n$ que $F$; si $e'_0$, $e''_0$ sont les images canoniques de $e'$, $e''$ dans $B_0$, ce sont deux idempotents orthogonaux tels que $1=e'_0+e''_0$, et $B_0$ est donc composée directe de $B'_0=e'_0B_0$ et $B''_0=e''_0B_0$. Soient $t$ et $t_0$ les images canoniques de $T$ dans $B$ et $B_0$; comme $B'_0$ (resp. $B''_0$) est une $k$-algèbre finie engendrée par $t'_0=e'_0t_0$ (resp. $t''_0=e''_0t_0$), elle admet une base de la forme $\{e'_0,t'_0,t_0'{}^2,\ldots,t_0'{}^{s-1}\}$ (resp. $\{e''_0,t''_0,t_0''{}^2,\ldots,t_0''{}^{r-1}\}$) avec $r+s=n$. D'autre part, $B$ étant un $A$-module libre (de base $\{1,t,\ldots,t^{n-1}\}$), $B'$ et $B''$ sont des $A$-modules projectifs, donc libres puisque $A$ est un anneau local (Bourbaki, \emph{Alg. comm.}, chap. II, \textsection~5, n\textsuperscript{o}~3, cor. de la prop. 5); il résulte donc de ce qui précède et de Bourbaki, \emph{Alg. comm.}, chap. II, \textsection~3, n\textsuperscript{o}~3, prop. 5, que si l'on pose $t'=e't$, $t''=e''t$, $\{e',t',\ldots,t'^{s-1}\}$ (resp. $\{e'',t'',\ldots,t''^{r-1}\}$) est une base du $A$-module $B'$ (resp. $B''$). Il y a donc un polynôme unitaire $H$ (resp. $G$) de degré $s$ (resp. $r$) de $A[T]$ tel que $e'H(t')=0$ et $e''G(t'')=0$; comme $t^h=t'^h+t''^h$ pour tout entier $h\geq1$, et $t'^h=e't'^h$, $t''^h=e''t''^h$, on a aussi $G(t)=e'G(t')$ et $H(t)=e''H(t'')$, d'où $G(t)H(t)=0$; on en conclut que le polynôme $G(T)H(T)$ est divisible par $F(T)$; mais comme les degrés de ces deux polynômes unitaires sont les mêmes, on a $GH=F$. En outre, $B'$ (resp. $B''$) est isomorphe à $A[T]/H.A[T]$ (resp. $A[T]/G.A[T]$). Enfin, il y a deux polynômes $R$, $S$ de $A[T]$ tels que $e'=R(t)$ et $e''=S(t)$; comme $R(t)=e'R(t')+e''R(t'')$, on a nécessairement $e''R(t'')=0$, et de même $e'S(t')=0$, de sorte que, par définition de $G$ et $H$, $R=QH$ et $S=PG$, où $P$, $Q$ appartiennent à $A[T]$; la relation $1=R(t)+S(t)$ dans $B$ donne donc par définition $PG+QH=1+LF$ pour un certain polynôme $L\in A[T]$, et comme $F=GH$, cela prouve que l'idéal engendré par $G$ et $H$ est $A[T]$, et achève de démontrer le lemme.

Ce lemme étant établi, il suffit de l'appliquer à l'anneau local $A$ d'une part, au corps $k$ de l'autre, pour voir aussitôt que les conditions $a')$ et $a'')$ sont équivalentes.

\begin{proposition}[18.5.14]
\label{IV.18.5.14}
Tout anneau semi-local $A$, séparé et complet pour la topologie $\mathfrak r$-préadique (où $\mathfrak r$ est le radical de $A$) est hensélien.
\end{proposition}

En effet, $A$ est composé direct d'anneaux locaux séparés et complets (Bourbaki, \emph{Alg. comm.}, chap. III, \textsection~2, n\textsuperscript{o}~13, cor. de la prop. 19), donc on est ramené au cas où $A$ est un anneau local. Vérifions le critère $a')$ de (18.5.11). Comme $B$ est un $A$-module libre de type fini, il est évidemment séparé et complet pour la topologie $\mathfrak r$-préadique,

